\scanpage{225}
% (9) and (10)
\[
 \sum D_i\omega^{\unsure{'}}_i=\overline\omega'\qquad\text{dans }\Phi
 \tag*{(9)}
\]
et la condition (7) signifie que l'on a, localement au dessus de $S$, une relation
\[
 \overline\omega\; \bigl(=\sum\overline D_{\deleted{\unsure{s}}i}\,\omega_i\bigr)=d\varphi,\qquad
 \varphi\in\Gamma(X,\mathcal O_X).\tag*{(10)}
\]

% period over a cycle
Pour voir que cette relation implique bien la nullité de (3), il suffit de prouver
\[
 \sum_i D_{i\,s}\int_{Z_s}\omega_{i\,s}=\int_{Z_s}\overline\omega
 \tag*{(11)}
\]
qui implique, si $\overline\omega=d\varphi$, que le premier mb est égal à $\int_{\partial Z_s}\varphi=0$ car $\partial Z_s=0$.

% general relation and its hypotheses
En fait, (11) provient de la relation générale
\[
 D_s\int_{Z_s}\omega_s=\int_{Z_s}\overline D_s\,\omega_s
 \tag*{(12)}
\]
valable pour tt op. diff. hol. $D$ sur $S$, et tte forme \inserted{holomorphe fermée} $\omega_{\deleted{\unsure{s}}}$ sur $X$ [le fait que $\omega$ est fermée implique que $\int_{Z_s}\omega$ est fonction holomorphe en $s$, et donne un sens au premier mb de (12)]. Pour prouver cette
% archive sheet number: 224
